% 01-intro.tex — 引言

\section{引言}
\label{sec:intro}

\subsection{问题: 元胞自动机生成的是迷宫吗?}

元胞自动机 (Cellular Automaton, CA) 是离散动力系统, 用局部更新规则在网格上同步演化。最早被研究的 Conway ``生命游戏'' 是 B/S 字符串规则: 8 邻域下, 活细胞在 $S$ 集合中邻居数时继续存活, 死细胞在 $B$ 集合中邻居数时复活 \cite{conway1970}. 这套规则在二维网格上能产生复杂图样, Wolfram 把它推广为四类动力学 \cite{wolfram2002}: 均匀, 周期, 混沌, 复杂 (生命游戏属于复杂类)。

CA 与迷宫生成在拓扑结构上很接近: 两者都关心二值网格, 都依赖邻域信息, 都讨论 ``局部规则涌现全局结构''。但 CA 的动力学是双向的, 可以把 ``看起来像迷宫的图样'' 视为 CA 的吸引子。这就引出一个具体的判定问题: 给定一个 CA 终态 (二值网格), 怎么判断它是否构成迷宫?

直觉上, 这个问题应该很简单 --- 数一下连通性、看一下走廊就行。但实际上, 至少有四类 CA 吸引子能骗过朴素判别:
\begin{itemize}[leftmargin=*, itemsep=0pt, topsep=0pt]
\item \textbf{实心块}. 墙比为 1, 整张图全黑, 拓扑完全退化。
\item \textbf{空盒}. 墙比为 0, 整张图全白, 同样退化。
\item \textbf{规则条纹 / 同心环 / 蜂巢}. 高度对称的几何反例, 看上去有结构, 但没有迷宫应有的一格宽走廊。
\item \textbf{随机噪声}. 50\% 密度的随机点能通过 ``连通性 $\ge$ 阈值'' 的检查, 但分叉过多, 不像迷宫。
\end{itemize}

要拒收这些反例, 需要一个同时刻画 ``走廊结构'' 和 ``非平凡拓扑'' 的度量, 还要在伪迷宫上判别失效。

\subsection{已有工作的两个空缺}

迷宫生成算法 (DFS \cite{pech2002}, Kruskal, Prim, Growing Tree, Sidewinder, Binary Tree) 与 CA 演化都能产生二值网格, 但两者之间缺一个统一的 ``是否迷宫'' 判别器。Bellot 等人 2021 年提出的 $F = \nu / \delta$ 度量是当前学界常用的对照基线, 但它对 ``少死端墙 + 长直径'' 的几何反例 (条纹、环、蜂巢) 反而给出高分 \cite{bellot2021}, 方向反了。

CA 规则搜索的现有工作集中在经典 B/S 字符串 ($2^{18}$ 候选), 用遗传编程或 ES 搜索 ``自然涌现迷宫'' 的规则 \cite{rejbrand2017, mcclendon2001}, 但搜索空间太小 (B/S 字符串只能产生有限种动力学), 不可能涌现复杂走廊结构。FamilyMask 编码把搜索空间扩大到 $2^{1648}$, 同时保留 B/S 字符串作为单族退化情形。

\subsection{本文贡献}

本文给出两个互补的产物:
\begin{enumerate}[leftmargin=*, itemsep=0pt, topsep=0pt]
  \item \textbf{FamilyMask 多族 CA 染色体。}把经典 B/S 规则推广为最多 16 族并行、按优先级仲裁的染色体。每族独立持有邻域 mask、出生集、存活集与优先级。搜索空间 $2^{103}$ 单族 / $2^{1648}$ 全染色体。
  \item \textbf{maze\_quality 八维几何度量。}5 拓扑 + 3 多样性子度量, 两层加权几何聚合 + $\min$ 平衡 + 墙比三角门控, 全部连续可微, ES 可学。在 15-pattern 基准上, 真迷宫均值 $0.711$, 伪迷宫均值 $0.000$, 间隔 $+0.711$, 误判 $0/9$。同数据集 Bellot 2021 $F$ 间隔 $-191.3$, 方向反转, 误判 $4/9$。
  \item \textbf{两个尺度的演化扫描。}小范围 ($40 \times 60$, 122 OK) 给出完整 (mask, maxFam) 消融数据, 最高分 $0.8233$ (manhattan-2 / maxFam=8 / seed 444)。大范围 ($500 \times 2000$, 5/5 OK) 验证 ``同配置, 大网格'' 的 scale-up, 最高分 $0.8195$, 5-run 均值 $0.8020$。两个 sweep 峰值差 $0.0037$, 表明小范围已接近该配置的上界。
\end{enumerate}

\subsection{全文结构}

§2 介绍相关工作 (CA 与迷宫、FamilyMask 类编码、几何度量、ES 搜索)。
§3 给出 FamilyMask 染色体的精确描述与执行语义。
§4 给出 maze\_quality 的 8 个子度量定义与两层聚合, 并在 15-pattern 基准上对比 Bellot 2021 $F$。
§5 报告两个 sweep 的实验设置、结果与对比。
§6 讨论两类失败模式 (预算受限 vs 表征受限) 与大网格 scale-up 的含义。
§7 总结。
附录 A 给出完整公式; 附录 B 给出数据可用性; 附录 C 给出软硬件环境。
